% ============================================================================
% ANEXO 1 --- MANUAL DO USUARIO
% ============================================================================

\section{Manual do Usu\'ario --- Guia Passo a Passo}
\label{sec:anexo1}

\subsection{O que \'e esta ferramenta?}

Imagine que voc\^e tem uma \textbf{tabela no Excel} com n\'umeros que representam o quanto cada poluente faz mal ao meio ambiente ou \`a sa\'ude das pessoas. Esses n\'umeros se chamam \textbf{fatores de caracteriza\c{c}\~ao}.

Agora imagine que existe um \textbf{programa de computador} chamado OpenLCA que os profissionais de meio ambiente usam para calcular os impactos de produtos. S\'o que esse programa n\~ao consegue ler sua tabela do Excel diretamente. Ele precisa receber os dados em um formato especial --- um ``pacote'' com muitos arquivos organizados de um jeito bem espec\'ifico.

\textbf{Esta ferramenta faz essa tradu\c{c}\~ao para voc\^e.} Ela pega sua tabela do Excel, organiza tudo no formato que o OpenLCA entende, e entrega um arquivo pronto para ser importado. Voc\^e n\~ao precisa saber programar.

\begin{lstlisting}[style=bash]
Sua tabela (Excel)  -->  [olca-cf-converter]  -->  Pacote pronto (ZIP)  -->  OpenLCA
\end{lstlisting}

\subsection{Gloss\'ario de termos t\'ecnicos}

\begin{longtable}{p{4cm}p{10cm}}
\toprule
\textbf{Termo} & \textbf{Significado em linguagem simples} \\
\midrule
\endhead
Fator de caracteriza\c{c}\~ao (CF) & Um n\'umero que diz ``quanto de impacto ambiental 1 kg de certo poluente causa''. Exemplo: 1 kg de am\^onia causa 0,000579 anos de vida perdidos (DALY). \\
\addlinespace
LCIA & Sigla para ``Avalia\c{c}\~ao de Impacto do Ciclo de Vida''. \'E a etapa em que os n\'umeros do invent\'ario s\~ao convertidos em impactos ambientais. \\
\addlinespace
OpenLCA & Um programa gratuito que profissionais usam para fazer Avalia\c{c}\~ao de Ciclo de Vida. Funciona como o ``Excel da ACV''. \\
\addlinespace
JSON-LD & O ``idioma'' que o OpenLCA fala. \'E um formato de arquivo. Voc\^e n\~ao precisa mexer nele --- a ferramenta cria automaticamente. \\
\addlinespace
UUID & Um c\'odigo \'unico (como um CPF para subst\^ancias qu\'imicas). O OpenLCA usa esses c\'odigos para conectar as coisas entre si. \\
\addlinespace
Ecoinvent & O maior banco de dados de invent\'ario ambiental do mundo. \\
\addlinespace
YAML & Um arquivo de texto simples onde voc\^e escreve configura\c{c}\~oes. Como preencher um formul\'ario em texto. \\
\addlinespace
Terminal & A tela onde voc\^e digita comandos. No Mac chama ``Terminal'', no Windows chama ``PowerShell''. \\
\addlinespace
Python & A linguagem de programa\c{c}\~ao em que a ferramenta foi escrita. Voc\^e precisa ter instalado, mas n\~ao precisa saber programar. \\
\addlinespace
Flow & Uma subst\^ancia emitida ao meio ambiente (ex: am\^onia, di\'oxido de enxofre). \\
\addlinespace
Endpoint & Dano final (ex: anos de vida perdidos --- DALY). \\
\addlinespace
Midpoint & Impacto intermedi\'ario (ex: kg de CO$_2$ equivalente). \\
\addlinespace
DALY & \textit{Disability-Adjusted Life Years} --- ``anos de vida saud\'avel perdidos''. Unidade que mede o dano \`a sa\'ude humana. \\
\bottomrule
\end{longtable}

\subsection{Pr\'e-requisitos}

Voc\^e precisa ter instalado no computador:

\begin{enumerate}
    \item \textbf{Python 3.10 ou mais recente} --- Para verificar, abra o Terminal e digite: \lstinline[style=bash]{python3 --version}
    \item \textbf{Acesso ao Terminal} --- Mac: aplicativo ``Terminal''; Windows: ``PowerShell''
\end{enumerate}

\subsection{Passo 1: Baixar a ferramenta}

Abra o Terminal e copie/cole estes comandos:

\begin{lstlisting}[style=bash]
git clone https://github.com/Roverlucas/Pos-Doc---ACV--OpenLCA-.git
cd Pos-Doc---ACV--OpenLCA-
\end{lstlisting}

O primeiro comando baixa os arquivos. O segundo entra na pasta que foi criada.

\subsection{Passo 2: Preparar o ambiente}

\begin{lstlisting}[style=bash]
python3 -m venv .venv
source .venv/bin/activate    # Mac/Linux
# .venv\Scripts\activate     # Windows
pip install -e ".[dev]"
\end{lstlisting}

Esses comandos criam um ``espa\c{c}o isolado'' e instalam a ferramenta. Voc\^e s\'o precisa fazer isso \textbf{uma vez}. Nas pr\'oximas vezes, basta ativar o ambiente com \lstinline[style=bash]{source .venv/bin/activate}.

\subsection{Passo 3: Preparar a planilha Excel}

Crie um arquivo \texttt{.xlsx} com 5 colunas:

\begin{table}[H]
\centering
\small
\begin{tabular}{lllll}
\toprule
\textbf{Flow} & \textbf{Category} & \textbf{Factor} & \textbf{Unit} & \textbf{Location} \\
\midrule
Ammonia & Elem. flows/Emission to air/unspec. & 0.000579 & DALY & Brazil \\
Nitrogen oxides & Elem. flows/Emission to air/high pop. & 0.001230 & DALY & Brazil, SP \\
\bottomrule
\end{tabular}
\end{table}

\textbf{O que cada coluna significa:}
\begin{itemize}
    \item \textbf{Flow} = Nome do poluente em ingl\^es (deve coincidir com o Ecoinvent)
    \item \textbf{Category} = ``Endere\c{c}o'' de onde a subst\^ancia \'e emitida
    \item \textbf{Factor} = O n\'umero do fator (com ponto decimal, n\~ao v\'irgula)
    \item \textbf{Unit} = Unidade do impacto (informativo)
    \item \textbf{Location} = Lugar para o qual o fator vale (``Global'' se n\~ao regionalizado)
\end{itemize}

\subsection{Passo 4: Criar a configura\c{c}\~ao}

\begin{lstlisting}[style=bash]
olca-cf init -o meu-metodo.yaml
\end{lstlisting}

Abra o arquivo \texttt{meu-metodo.yaml} em qualquer editor de texto e preencha com os dados do seu m\'etodo:

\begin{lstlisting}[style=yaml]
method:
  name: "Meu Metodo LCIA"
  description: "Descricao com referencia bibliografica."
category:
  name: "minha categoria de impacto"
  type: endpoint
units:
  input:
    name: kg
    property: Mass
  output:
    name: DALY
    property: "Impact on human health"
files:
  excel: "meus-fatores.xlsx"
  reference_zip: "Base-Ecoinvent.zip"
  output_zip: "Meu-Metodo-FINAL.zip"
\end{lstlisting}

\subsection{Passo 5: Verificar a planilha (recomendado)}

\begin{lstlisting}[style=bash]
olca-cf validate meus-fatores.xlsx
\end{lstlisting}

Se tudo estiver correto, aparecer\'a: ``Excel file is valid!''. Se houver erro, a ferramenta dir\'a exatamente o que precisa ser corrigido.

\subsection{Passo 6: Converter}

\begin{lstlisting}[style=bash]
olca-cf convert meu-metodo.yaml
\end{lstlisting}

O arquivo \texttt{Meu-Metodo-FINAL.zip} ser\'a criado --- este \'e o pacote para importar no OpenLCA.

\subsection{Passo 7: Importar no OpenLCA}

\begin{enumerate}
    \item Abra o OpenLCA
    \item Menu \textbf{File} $\rightarrow$ \textbf{Import}
    \item Selecione \textbf{Linked Data (JSON-LD)}
    \item Navegue at\'e o arquivo \texttt{.zip} gerado
    \item Clique em \textbf{Finish}
\end{enumerate}

O m\'etodo aparecer\'a em \textbf{Impact assessment methods} no navegador.

\subsection{Resumo r\'apido}

\begin{lstlisting}[style=bash]
# Instalacao (uma vez):
git clone https://github.com/Roverlucas/Pos-Doc---ACV--OpenLCA-.git
cd Pos-Doc---ACV--OpenLCA-
python3 -m venv .venv && source .venv/bin/activate
pip install -e ".[dev]"

# Uso (toda vez):
source .venv/bin/activate
olca-cf validate meu-excel.xlsx
olca-cf convert meu-config.yaml
\end{lstlisting}

\subsection{Perguntas frequentes}

\begin{description}
    \item[Preciso saber programar?] N\~ao. Voc\^e s\'o prepara a planilha e o arquivo de configura\c{c}\~ao.
    \item[Funciona no Windows?] Sim. Funciona em Windows, Mac e Linux.
    \item[E se eu n\~ao tenho a base do Ecoinvent?] Funciona sem ela. Coloque \texttt{reference\_zip: null} no config.
    \item[Posso usar para qualquer tipo de impacto?] Sim. A ferramenta \'e gen\'erica.
    \item[Meus fatores n\~ao s\~ao regionalizados.] Coloque ``Global'' na coluna Location.
\end{description}
